% Source pp. 161-172 (Lectures13-18.pdf).
\section{Wilson's Renormalization Group}

\subsection*{On Second Order Phase Transitions}

Let us briefly discuss the so called \emph{continuous} or \emph{second order}
phase transitions, because it is the theory of continuous phase transitions in
statistical mechanics which is related to (Euclidean) quantum field theory in
most close way.

Consider a magnet, i.e.\ a piece of a magnetic material. It can be thought of
as consisting of ``micromagnets'', i.e.\ small elements having nonzero magnetic
moments. Although in reality the structure of ``micromagnets'' is usually rather
complex, for the sake of this discussion we can think of, say, a cubic lattice
with each cell representing a ``micromagnet'' having a magnetic moment
(proportional to the volume of the cell),
\[
\diac[0.55]{\draw[prop] (0,0) rectangle (2,2); \draw[prop] (0,2) -- (0.8,2.8) -- (2.8,2.8) -- (2,2); \draw[prop] (2.8,2.8) -- (2.8,0.8) -- (2,0);
  \draw[densely dashed] (0,0) -- (0.8,0.8) -- (2.8,0.8) (0.8,0.8) -- (0.8,2.8);
  \draw[{Stealth[length=1.4mm]}-{Stealth[length=1.4mm]}] (3.2,0.8) -- node[right] {\small$a$} (3.2,2.8);}
\qquad\longrightarrow\qquad a^3\cdot\mathbf m_i
\]
The interaction between the ``micromagnets'' is such that they tend to align,
thus producing an overall magnetic moment
\[
V\,\mathbf M=V\,\langle\mathbf m_i\rangle\,,
\]
where $\mathbf M$ denotes specific (per unit volume) magnetization, and the
expectation value in the r.h.s.\ is taken over thermal equilibrium state. The
thermal fluctuations tend to put the directions of $\mathbf m_i$ at random, and
as the result of this competition the ferromagnet has two phases, the low
temperature, or \emph{ordered} phase with nonvanishing local magnetization
$\mathbf M$, and high-temperature \emph{disordered} phase with $\mathbf M=0$.

The most important qualitative features of this phase transition, which occurs
at certain critical temperature
\[
T=T_c\,,
\]
can be understood through the \emph{Landau Theory}.

\paragraph{Landau Theory:} Consider the thermodynamic potential
\[
V\,\Omega(\mathbf M,T)
\]
as the function of specific magnetization $\mathbf M$ and temperature $T$;
$\Omega$ denotes the \emph{specific} (per unit volume) thermodynamic potential.
The equilibrium value of $\mathbf M$ is determined by the equation
\[
\frac{\partial\Omega(\mathbf M,T)}{\partial\mathbf M}=0\,.
\]
Consider the high-temperature phase $T>T_c$, where the solution of this equation
is $\mathbf M=0$, and assume that one can expand $\Omega(\mathbf M,T)$ in powers
of $\mathbf M$ near this minimum,
\[
\Omega(\mathbf M,T)=A(T)+B(T)\,\mathbf M^2+C(T)\,\mathbf M^4+\cdots
\]
(I assume here isotropy of the problem, hence the expansion is in powers of
$\mathbf M^2$). It is clear that $\mathbf M=0$ is the right solution (i.e.\ the
minimum of $\Omega$) only as long as
\[
B(T)>0\ \to\ \mathbf M=0\,.
\]
When $B(T)$ becomes negative (as it happens at $T<T_c$, according to this
theory), the minimum occurs at nonzero $\mathbf M$ such that
\[
|\mathbf M|=\sqrt{\frac{-B}{2C}}\,;
\]
in fact there is a continuum of degenerate minima corresponding to different
directions of $\mathbf M$. This situation represents the low-temperature phase
where spontaneous magnetization appears. Thus, in the Landau theory the critical
temperature $T_c$ is the point where $B(T)$ changes its sign,
\[
\begin{array}{lll}
B(T)>0 & \text{at} & T>T_c\\
B(T)<0 & \text{at} & T<T_c
\end{array}
\]
Landau has assumed that near $T_c$ one can write $B(T)$ as
\[
B(T)=b\,(T-T_c)+\cdots\,,
\]
with some positive constant $b$. Then it follows that
\[
\begin{array}{lll}
|\mathbf M|=0 & \text{at} & T>T_c\\
|\mathbf M|\sim(T_c-T)^{1/2} & \text{at} & T<T_c
\end{array}
\]
It is not difficult to calculate the entropy at $T\approx T_c$,
\[
S/V=-\frac{\partial\Omega(\mathbf M,T)}{\partial T}\approx S_0/V+
\begin{cases}0 & \text{for}\quad T>T_c\\ \frac{b^2}{2C}\,(T-T_c) & \text{for}\quad T<T_c\end{cases}
\]
where $S_0$ refers to the critical point. The entropy (and hence the internal
energy) is continuous at $T_c$, so there is no latent heat. This is the
characteristic property of a second order phase transition. The specific heat
\[
C=T\,\frac{\partial S}{\partial T}
\]
has finite discontinuity at $T=T_c$.

\paragraph{Critical behavior:} Although the Landau theory explains existence and
the ``continuous'' nature of the critical point, in real systems the singularity
of the thermodynamic functions at $T_c$ turns out to be more complicated than the
above arguments predict. Experiment shows that
\[
|\mathbf M|\sim(T_c-T)^\beta\qquad\text{for}\quad T<T_c\,,
\]
with $\beta\approx0.35\pm0.03$, definitely different from $0.5$. The specific heat
is not only discontinuous, but in fact diverges near $T_c$,
\[
C\sim|T_c-T|^{-\alpha}\,,
\]
with $\alpha\approx0.12$, not 0.

The Landau theory does not take into account prominent role of the
\emph{long-range fluctuations} of the magnetization near critical point.
According to the thermodynamic theory of fluctuations, the probability to have
some non-equilibrium value $\mathbf M$ of the magnetization is (for $T>T_c$)
\[
P(\mathbf M)\sim e^{-\frac{V\Omega(\mathbf M,T)}{kT}}\sim e^{-\frac{VB(T)}{kT}\,\mathbf M^2}
\]
If $B>0$, in the limit $V\to\infty$ the fluctuations are totally suppressed.
Vanishing of $B(T)$ at $T=T_c$ indicates that in the critical domain the
fluctuations must play significant role.

\subsection*{Fluctuation Theory of Phase Transitions}

Previously, we have given a brief discussion of critical behavior near a second
order phase transition. We have pointed out singular character of the
thermodynamic functions near the critical point, and suggested that the origin of
the singularities is in the enhanced role of fluctuations.

The qualitative picture of the fluctuations is as follows. In the vicinity of
$T_c$ the competition between the energetic tendency of the micromagnets to
align, and the thermal fluctuations leads to formation of large domains of
approximately collinear micromagnets (figure~\ref{fig:domains}),
\begin{figure}[ht]
\centering\figfile{domains}
\caption{Two distant domains near $T_c$. Inside each domain the micromagnets are
nearly aligned; the average directions in different domains differ.}
\label{fig:domains}
\end{figure}
so that if one looks inside one domain the magnets are almost aligned, but their
average directions can be largely different in two distant domains. Near $T_c$
the average sizes of the domains become much larger than any microscopic scales,
and the domains involve macroscopic numbers of the micromagnets. The orientations
of the domains experience fluctuations which vary slowly in space (in the
microscopic scale). These long-range fluctuations lead to formation of the
critical singularities at $T_c$.

Having this picture in mind, one can consider local magnetization
$\mathbf M(x)$, a slow varying function of the spatial coordinates (meaning that
the distance at which $\mathbf M(x)$ changes significantly must involve
macroscopic number of the micromagnets). One then replaces the homogeneous
magnetization $\mathbf M$ by the local magnetization
\[
V\,\Omega(\mathbf M,T)\ \to\ \int d^3x\,\Omega\bigl(\mathbf M(x),T\bigr)\,.
\]
Once $\mathbf M$ is not a constant, one may include the spatial derivatives of
$\mathbf M(x)$, and assume that the long-range fluctuations have the probability
distribution
\[
P[\mathbf M(x)]\sim Z^{-1}\exp\left(-\frac{\Ham[\mathbf M(x)]}{kT}\right)
\]
with the ``effective energy'' of the form
\[
\Ham[\mathbf M(x)]=\int d^3x\,\Omega\bigl(\mathbf M(x),\nabla\mathbf M(x),\dots;T\bigr)
\]
where $\Omega$ now is a local density. In the spirit of the Landau theory we can
write
\[
\Omega=A(T)+B(T)\,\mathbf M^2(x)+C(T)\,\mathbf M^4(x)+D(T)\,\bigl(\nabla\mathbf M(x)\bigr)^2+\dots
\]
where the dots may involve higher powers of $\mathbf M^2$ as well as terms with
higher derivatives. Neglecting all but the terms explicitly shown above
corresponds to the so-called Landau-Ginzburg model.

With this probability distribution, the free energy of the magnet is to be
computed as some sort of functional integral
\[
Z=e^{-\frac F{kT}}=\int\bigl[D\mathbf M(x)\bigr]\,e^{-\frac1{kT}\Ham[\mathbf M(x)]}
\]
where the integration is over the functions $\mathbf M(x)$ which vary slowly in
terms of the microscopic scale.

Note that this expression differs mostly in notations from the functional
integral in Euclidean scalar field theory. After the rescaling
\[
\sqrt{\frac{2D(T)}{kT}}\,\mathbf M(x)=\varphi_0(x)
\]
it becomes
\[
Z=\int\bigl[D\varphi_0\bigr]\,e^{-\mathcal A[\varphi_0]}\,,
\]
where in the Landau-Ginzburg approximation
\[
\mathcal A=\int\left[\frac12\,(\partial\varphi_0)^2+\frac{m_0^2}2\,\varphi_0^2+\frac{\lambda_0}{4!}\,\varphi_0^4\right]
\]
has the familiar form, with
\[
m_0^2=m_0^2(T)\,,\qquad\lambda_0=\lambda_0(T)
\]
some functions of $T$, whose form depends on microscopic details of the magnetic
system. On the superficial level the main difference is in the dimensionality of
the space (3 instead of 4), and the number of components of the scalar field (3
instead of 1). As for the difference in $d$, we have already seen that the idea
of studying QFT in spaces of different (even non-integer) dimensions $d$ is
useful. The problematic of the phase transition theory, in its application to
condensed matter physics, usually corresponds to the case $d=3$. But it is
interesting and instructive to study criticality at any $d$, in particular to see
how the characteristics of a phase transition change with $d$. Also, in the above
simple model of magnet, the direction of $\mathbf M$ is completely decoupled from
any space direction. Therefore formally the dimensionality of the vector
$\mathbf M$, or equivalently the number of components of
\[
\varphi_0=\{\varphi_{0,i}\}\,,\qquad i=1,2,\dots,N
\]
can be any number, not necessarily $d$. In fact, there are magnets with
anisotropy of the type of ``easy plane'' or ``easy axis'', which correspond to
$N=2$ and $N=1$, respectively. In most of the following discussion I will assume
$N=1$ for simplicity of notations.

On a deeper level, one can raise two objections against this picture. First, the
above formulation assumes integration only over the ``slow varying'' fields
$\varphi_0(x)$. This is similar to the idea of cutoff in field theory, but here
it seems to be dictated by physics of the problem. Indeed, in our approach the
``thermodynamic potential''
\[
\Ham([\mathbf M(x)],T)
\]
already takes into account ``microscopic'' thermal fluctuations in our system,
while the ``macroscopic'' fluctuations are left to be accounted for in the
functional integral over $D[\mathbf M(x)]$. But how these two types of
fluctuations are separated? We have said that $\mathbf M(x)$ can vary on the
scale ``much larger than the microscopic scales''. But how much larger? Should we
take the cutoff parameter $\Lambda=(100\,a)^{-1}$, or $(10^4\,a)^{-1}$? The second
objection is that it is not clear what makes it possible to neglect the higher
order (both in $\mathbf M$ and derivatives) terms in $\Ham$.

These questions are at the heart of the Wilson's approach to the Renormalization
Group.

\subsection*{Wilson's Approach\protect\footnote{Discussion in this and following
Sections has overlap with \S12.1 in PS. More systematic treatment is in K.~Wilson,
J.~Kogut, ``The RG and epsilon expansion'', Phys.\ Rep.\ \textbf{12C} (1974)
(highly recommended!). See also L.~Kadanoff, ``Statistics, Dynamics, and
Renormalization'', WS.}}

To come to the idea of this approach let us ask instead \emph{what difference
does it make} if we take
\[
\Lambda=\Lambda_0=(100\,a)^{-1}\qquad\text{or}\qquad\Lambda=\Lambda_1=(10^4\,a)^{-1}\,.
\]
In other words, if we consider say $\varphi^4$ theory, and write down the
functional integral
\[
\int D[\varphi_0]\,e^{-\mathcal A[\varphi_0]}\,(\dots)
\]
in its bare form, with some cutoff regularization implied, how the theory depends
on the value of the cutoff parameter $\Lambda$?

Because basic ideas of the Wilson's RG do not depend very much on precise way how
the cutoff regularization is implemented, let me assume specific one. Let me
assume that the field $\varphi_0$ can be written as the Fourier-like integral
\[
\varphi_0(x)=\int_{|k|<\Lambda}\frac{d^dk}{(2\pi)^d}\,\phi(k)\,e^{ikx}\,,
\]
where only the wave-vectors with $|k|<\Lambda$ contribute to the integral. This
can be named (as I will) the ``Sharp momentum cutoff'', because in perturbation
theory it corresponds to the cutoff propagator
\[
D_\Lambda(k)=\frac{\Phi(k^2/\Lambda^2)}{k^2+m_0^2}
\]
with $\Phi(x)=\Theta(1-x)$ being the step function. (We did not consider this
regularization before because it is not very convenient in evaluating usual
Feynman diagrams).

So our question now is how the functional integral depends on $\Lambda$, or what
is the difference between the integral with $\Lambda=\Lambda_0$ and
$\Lambda=\Lambda_1$. Let us temporarily use the notations
\[
\varphi_0(x)=\int_{|k|<\Lambda_0}\phi(k)\,e^{ikx}\,,\qquad
\varphi_1(x)=\int_{|k|<\Lambda_1}\phi(k)\,e^{ikx}\,.
\]
If $\Lambda_0>\Lambda_1$ the difference of course is that the integration over
$\varphi_0$ includes the integration over all modes with the momenta
\[
\Lambda_0>|k|>\Lambda_1\,,
\]
while the integral over $\varphi_1$ does not. We can write
\[
\varphi_0(x)=\varphi_1(x)+\tilde\varphi(x)\,,\qquad\text{where}\qquad
\tilde\varphi(x)=\int_{\Lambda_0>|k|>\Lambda_1}\phi(k)\,e^{ikx}\,.
\]
In this context, the parts $\varphi_1(x)$ and $\tilde\varphi$ are often called
the ``slow part'' and the ``fast part'', respectively. We have
\[
D[\varphi_0]=D[\varphi_1]\,D[\tilde\varphi]\,.
\]
Suppose we want to calculate some correlation function in the theory with the
cutoff $\Lambda_0$,
\[
\bigl\langle\varphi_0(x_1)\cdots\varphi_0(x_n)\bigr\rangle=\frac1{Z_0}\int D[\varphi_0]\,\varphi_0(x_1)\cdots\varphi_0(x_n)\,e^{-\mathcal A_0[\varphi_0]}\,,
\]
where $\mathcal A_0$ is some action, say $\varphi^4$ action above. We are
interested in low-momentum behaviour of the associated proper vertices, with the
momenta conjugated to $x_1,\dots,x_n$ much smaller than both $\Lambda_0$ and
$\Lambda_1$. Therefore in the insertions we can replace
\[
\varphi_0(x_1)\cdots\varphi_0(x_n)\to\varphi_1(x_1)\cdots\varphi_1(x_n)\,.
\]
Then the functional integral for this correlation function becomes
\[
\frac1{Z_0}\int D[\varphi_1]\,\varphi_1(x_1)\cdots\varphi_1(x_n)\int D[\tilde\varphi]\,e^{-\mathcal A_0[\varphi_1+\tilde\varphi]}\,.
\]
Now, the idea is to integrate over the fast part $\tilde\varphi$ first. This
operation defines the new ``effective action'' $\mathcal A_1[\varphi_1]$ for the
slow part,
\[
\int D[\tilde\varphi]\,e^{-\mathcal A_0[\varphi_1+\tilde\varphi]}=\Const\,e^{-\mathcal A_1[\varphi_1]}\,.
\]
The constant factor here is not important, because after the same operation it
appears in $Z_0$ as well. Thus, we have
\[
\text{f.t.}\,\bigl\langle\varphi_0(x_1)\cdots\varphi_0(x_n)\bigr\rangle\Big|_{\mathcal A_0,\Lambda_0}
=\text{f.t.}\,\bigl\langle\varphi_1(x_1)\cdots\varphi_1(x_n)\bigr\rangle\Big|_{\mathcal A_1,\Lambda_1}\,,\qquad p_i\ll\Lambda\,.
\]
In other words, if only the low-momentum correlation functions are of concern,
the theory with cutoff $\Lambda_0$ and action $\mathcal A_0$ is equivalent to the
theory with lower cutoff $\Lambda_1$ and new action $\mathcal A_1$,
\[
\mathcal A_0,\Lambda_0\ \to\ \mathcal A_1,\Lambda_1\,.
\]
Note that I am not writing $\leftrightarrow$ here, because this operation can not
be undone. Information about the short-distance interactions, of the scales
between $\Lambda_0^{-1}$ and $\Lambda_1^{-1}$ is lost.

The new action $\mathcal A_1$ is defined through the functional integral over the
fast part $\tilde\varphi$. Although at this level of discussion precise form of
$\mathcal A_0$ is not important, to get some idea what this operation amounts to,
let us assume for the moment that $\mathcal A_0$ is the $\varphi^4$ action.
Substituting
\[
\varphi_0=\varphi_1+\tilde\varphi
\]
we find
\[
\mathcal A_0[\varphi_0]=\mathcal A_0[\varphi_1]+\tilde{\mathcal A}[\tilde\varphi,\varphi_1]\,,
\]
and after the integration over $\tilde\varphi$
\[
\mathcal A_1[\varphi_1]=\mathcal A_0[\varphi_1]+\mathcal A'[\varphi_1]\,,\qquad\text{where}\qquad
\Const\,e^{-\mathcal A'[\varphi_1]}=\int[D\tilde\varphi]\,e^{-\tilde{\mathcal A}[\tilde\varphi,\varphi_1]}\,.
\]
For $\mathcal A_0$ of the $\varphi^4$ form
\[
\tilde{\mathcal A}=\int d^dx\,\left[\frac12\,(\partial\tilde\varphi)^2+\frac{m_0^2}2\,\tilde\varphi^2
+\frac{\lambda_0}{4!}\,\bigl(4\varphi_1^3\tilde\varphi+6\varphi_1^2\tilde\varphi^2+4\varphi_1\tilde\varphi^3+\tilde\varphi^4\bigr)\right]
\]
The terms $\partial\varphi_1\partial\tilde\varphi$ and $\varphi_1\tilde\varphi$ do
not appear because by definition $\varphi_1$ and $\tilde\varphi$ carry different
momenta.

The integral over $\tilde\varphi$ is complicated non-gaussian integral. The
functional $\tilde{\mathcal A}[\tilde\varphi,\varphi_1]$ in this context should be
regarded as the action for $\tilde\varphi$, with $\varphi_1$ determining the
``source terms'' (linear, quadratic and cubic in $\tilde\varphi$).

As usual, the functional integral over $\tilde\varphi$ can be represented as a sum
of the Feynman diagrams. Let us denote the $\tilde\varphi$-propagator by a solid
line,
\[
\langle\tilde\varphi\tilde\varphi\rangle\ \to\
\dia{\draw[line width=0.9pt] (-1.6,0) -- (1.6,0); \node[above] at (0,0.05) {\small$k$};}
=\frac{\Theta(\Lambda_0>|k|>\Lambda_1)}{k^2+m_0^2}\,,
\]
with obvious meaning of the step-function $\Theta$. The remaining terms in
$\tilde{\mathcal A}$ give rise to the vertices
\begin{align*}
\dia[0.45]{\foreach \a in {150,180,210} \draw[prop,densely dashed] (0,0) -- (\a:1.4); \draw[line width=0.9pt] (0,0) -- (1.6,0); \node[vtx] at (0,0) {};}
&=-\frac{\lambda_0}{3!}\int_x\varphi_1^3(x)\,,&
\dia[0.45]{\foreach \a in {135,225} \draw[prop,densely dashed] (0,0) -- (\a:1.4); \foreach \a in {45,315} \draw[line width=0.9pt] (0,0) -- (\a:1.4); \node[vtx] at (0,0) {};}
&=-\frac{\lambda_0}2\int_x\varphi_1^2(x)\,,\\
\dia[0.45]{\draw[prop,densely dashed] (-1.4,0) -- (0,0); \foreach \a in {45,0,-45} \draw[line width=0.9pt] (0,0) -- (\a:1.4); \node[vtx] at (0,0) {};}
&=-\lambda_0\int_x\varphi_1(x)\,,&
\dia[0.45]{\foreach \a in {45,135,225,315} \draw[line width=0.9pt] (0,0) -- (\a:1.4); \node[vtx] at (0,0) {};}
&=-\lambda_0\int_x\,.
\end{align*}
Here the dashed line represents the field $\varphi_1$, the number of the dashed
lines showing the power of $\varphi_1$. Each dashed line carries momentum with
$|k|<\Lambda_1$.

We have to integrate over $\tilde\varphi$, with $\varphi_1$ fixed. Therefore all
the solid lines in the diagrams are contracted, while all the dashed lines are
left as external legs.

As usual, the sum of all the diagrams exponentiate in terms of the sum of
connected diagrams, i.e.\ the functional $\mathcal A'[\varphi_1]$ is given by the
sum of all connected diagrams,
\[
\mathcal A'[\varphi_1]=-\bigl(\text{sum of all connected diagrams}\bigr)
\]
Also, the diagrams with no external legs can be ignored, because they only
contribute to the constant factor in $\exp(-\mathcal A')$.

It is obvious that the new action $\mathcal A_1=\mathcal A_0+\mathcal A'$ does not
have simple form of the $\varphi^4$ action. The contributions to $\mathcal A'$
contain all powers of $\varphi_1$, from diagrams like this
\[
\diac[0.45]{\draw[line width=0.9pt] (0,0) circle (1); \coordinate (a) at (180:1); \coordinate (b) at (50:1); \coordinate (c) at (-60:1); \coordinate (d) at (-115:1);
  \draw[line width=0.9pt] (d) .. controls (0,-0.2) .. (c);
  \foreach \p/\q in {a/135,a/225,b/90,b/20,c/-20,d/-110} \draw[prop,densely dashed] (\p) -- ++(\q:1.3);
  \foreach \p in {a,b,c,d} \node[vtx] at (\p) {};}
\ \sim\varphi^6
\]
Moreover, the action $\mathcal A_1$ is not even formally local. For example, the
diagram
\[
\diac[0.45]{\draw[line width=0.9pt] (0,0) ellipse (1 and 0.7); \coordinate (L) at (-1,0); \coordinate (R) at (1,0);
  \foreach \p/\q in {L/135,L/225,R/45,R/315} \draw[prop,densely dashed] (\p) -- ++(\q:1.4);
  \node[vtx] at (L) {}; \node[vtx] at (R) {};}
\]
brings in the contribution
\[
-\frac{\lambda_0^2}{16}\int d^dx\,d^dx'\,\varphi_1^2(x)\,\tilde D^2(x-x')\,\varphi_1^2(x')\,,
\]
where $\tilde D$ is the $\tilde\varphi$ propagator,
\[
\tilde D(x)=\int_{\Lambda_0>|k|>\Lambda_1}\frac{d^dk}{(2\pi)^d}\,\frac{e^{ikx}}{k^2+m_0^2}\,.
\]
This non-locality, albeit unpleasant, is not disastrous. Recall that we have
already violated the locality by introducing the cutoff. It is important however
that the additional nonlocality in $\mathcal A_1$ is of short-range nature. More
precise statement is that the new non-local terms generated in this way admit
convergent derivative expansion.

This means the following. Suppose we have some diagram evaluated in the momentum
space,
\[
\diac[0.45]{\draw[line width=0.9pt] (0,0) circle (0.8); \node[font=\tiny] at (0,0) {diagram};
  \foreach \a/\l in {135/p_n,90/p_1,45/p_2,0/p_3} {\draw[prop] (\a:0.8) -- (\a:2.2); \node[vtx] at (\a:0.8) {}; \node at (\a:2.6) {\small$\l$};}
  \draw[densely dotted] (165:1.5) arc[start angle=165,end angle=330,radius=1.5];}
=\mathcal D(p_1,\dots,p_n)\,.
\]
One can expand $\mathcal D(p_1,\dots,p_n)$ in a simultaneous power series in all
the external momenta $p_i$, and the series will converge all the way up to the
momenta $|p_i|\le\Lambda_1$. This convergence is the consequence of the fact that
all the integration momenta in the diagram lie above $\Lambda_1$. Note that
should the integration involve all momenta, as in usual Feynman diagrams, similar
expansion would converge only up to $|p_i|<m_0$, at the best.

This convergence of the expansions of $\mathcal D(p_1,\dots,p_n)$ implies that the
new action $\mathcal A_1$ admits meaningful expansion in terms of local
composites,
\[
\mathcal A_1[\varphi_1]=\int d^dx\,\left[\sum_{k=1}^\infty\frac{u_{2k}}{2k!}\,\varphi_1^{2k}(x)+\sum_{k=0}^\infty\frac{v_{2k}}{2k!}\,\varphi_1^{2k}(x)\,\bigl(\partial\varphi_1(x)\bigr)^2+\text{higher terms}\right]
\]
where the ``higher terms'' involve both the higher powers of $\partial\varphi_1$,
and the terms with the higher derivatives of $\varphi_1$. The most significant
difference from truly local action is that these higher terms include derivatives
of arbitrarily high order. However it is assumed that the full expansion
converges for the fields $\varphi_1$, in which only the Fourier modes with
$|k|<\Lambda_1$ are present. I will call this form ``quasi-local''.

Our conclusion at this point is that by integrating out the fast part
$\tilde\varphi$, the theory described originally by the action $\mathcal A_0$,
with the cutoff $\Lambda_0$, can be equivalently described by new action
$\mathcal A_1$, with the lower cutoff $\Lambda_1$,
\[
\mathcal A_0,\Lambda_0\ \to\ \mathcal A_1,\Lambda_1\,.
\]
The new action however has complicated quasi-local form.

At first glance it may look as a step in wrong direction. But it becomes very
powerful move when combined with another step, which is to make a variable
transformation
\[
\varphi_1(x)=Z^{-1/2}(L)\,\varphi'_0(x/L)
\]
in the functional integral
\[
\frac1{Z_1}\int[D\varphi_1]\,\varphi_1(x_1)\cdots\varphi_1(x_n)\,e^{-\mathcal A_1[\varphi_1]}\,,
\]
where $L$ is the ratio
\[
L=\Lambda_0/\Lambda_1>1\,,
\]
and $Z^{1/2}(L)$ is a factor to be discussed later. This is the dilation with the
factor $1/L$, which returns the cutoff from $\Lambda_1$ back to the original value
$\Lambda_0$. We have
\[
\varphi_1(x)=\int_{|k|<\Lambda_1}\phi(k)\,e^{ikx}=\int_{|k|<\Lambda_1}\phi(k)\,e^{i(Lk)(x/L)}=L^{-d}\int_{|q|<\Lambda_0}\phi(q/L)\,e^{iq(x/L)}\,,
\]
so that
\[
\varphi'_0(x)=L^{-d}\,Z^{1/2}(L)\int_{|k|<\Lambda_0}\phi(k/L)\,e^{ikx}
\]
indeed contains all the wave-vectors with $|k|<\Lambda_0$.

For the action this amounts to transformation
\[
\mathcal A_1[\varphi_1]=\mathcal A'_0[\varphi'_0]\,,
\]
where $\mathcal A'_0$ is obtained from $\mathcal A_1$ simply by re-scaling of the
coordinates, plus the field renormalization. We then have for correlation
functions
\[
\bigl\langle\varphi_0(x_1)\cdots\varphi_0(x_n)\bigr\rangle\Big|_{\mathcal A_0,\Lambda_0}
=Z^{-n/2}(L)\,\bigl\langle\varphi_0(x_1/L)\cdots\varphi_0(x_n/L)\bigr\rangle\Big|_{\mathcal A'_0,\Lambda_0}\,,
\]
i.e.\ the theories described by $\mathcal A_0$ and $\mathcal A'_0$, \emph{with the
same cutoff} $\Lambda_0$ are different only by the scale transformation, with the
scale factor $L$. The transformation
\[
RG(L):\quad\mathcal A_0\to\mathcal A'_0
\]
is the Wilson's RG transformation. It represents the scale transformation in the
space of effective actions.
